\section{Introduction}
\label{sec:intro}

Text-to-video models such as Wan~\cite{wan2025wan}, CogVideoX~\cite{yang2025cogvideox} and VACE~\cite{jiang2025vace} synthesise convincing street scenes, shop fronts and screens, yet the text inside them is often garbled: a record shop reads \str{ROCRDS}, a convenience store reads \cjkimg{cjk_eab75b7f4d} (\cref{fig:teaser}). In our benchmark of generated clips, only \rOrigRead\% of the target strings are rendered correctly. Wherever the text carries meaning, a single wrong character is a hard failure, and re-generating the clip until the text comes out right is expensive and changes the shot. We therefore study \emph{repair}: keep the generated clip, and make its text read correctly.

The natural tool is a generative video text editor, which regenerates the text region conditioned on rendered glyphs and the source style. On generated videos we observe three systematic problems. Such editors write characters that look plausible but are wrong---\str{FRESY BREAD}, \zh{小心坍滑}---and the strongest of them, \steer~\cite{zeng2026steervte}, still leaves \rSteerUnread\% of our clips unreadable. They rewrite text that was already correct (\str{FRESH BREAD} $\to$ \str{FRECH BREAD}, \cref{fig:teaser}c). And because the whole clip passes through the generator, the background changes too: PSNR outside the text region drops to \rSteerBG\,dB.

We argue that repair should be correctness-first and minimal: do not change what is already right; when a change is needed, change only the text; and recover the material of the old text only when there is evidence that this helps. \method implements this order with three evidence gates (\cref{fig:overview}). A preservation gate returns the clip unchanged when OCR on many frames already reads the target. A localization gate decides whether the new text can be laid out on the old glyph instance---its size, colour and baseline---or must fall back to filling the tracked box. An appearance gate decides whether to apply source appearance transport (SAT), a non-parametric operator that carries glow and material from the old glyphs to the new ones. Between the gates the target is rendered deterministically from a matched font and warped with carrier homographies, so its spelling cannot drift, and only pixels inside the text region are composited. Nothing is trained.

On 526 generated clips covering seven carrier materials in English and Chinese, \method raises readability from \rOrigRead\% to \rOursRead\% (\steer: \rSteerRead\%), with \rOursNet{} net repairs against \rSteerNet, and keeps the background intact (\rOursBG\,dB). We are equally explicit about what it does not do: its glyphs are rendered, not generated, and where both methods repair a clip, \steer's text is closer to the original appearance in \rSteerCloserPct\% of the cases. We trace this gap to its causes and report the alternatives we tried and rejected under pre-registered criteria.

Our contributions are: \textbf{(i)} a repair task, a 526-clip benchmark and a protocol whose metrics separate repairs from breakages, measure background fidelity and temporal stability, and are complemented by blind human studies and a pre-registered second reader; \textbf{(ii)} \method, a training-free, evidence-gated pipeline with deterministic rendering anchored to the tracked carrier and SAT behind an appearance gate; \textbf{(iii)} an analysis with negative results---a learned 1.3B refiner, two appearance modules, a glyph-locked material generator and a carrier-grounded branch for inputs without any source glyph, each judged by criteria fixed before the run---and a verified selection between \method and \steer that reaches \hVRead\% readability with \hVBroken{} broken clips (\cref{sec:exp_hybrid}). The full experimental record, including every pre-registration and all audits, is in the supplementary material.
